\section{相关工作与贡献边界}\label{sec:related}
\paragraph{关系抽象与 SMT。}
Haas 等~\cite{haas2023} 对内存模型关系建立 may/must 上下界，利用过滤器和外部语义信息简化 SMT 编码。这是本文最接近的分析对象。$R_\forall$ 可以作为其语义过滤器，第~\ref{sec:native}节还给出实际接入结果。因此本文不声称首次发现“不确定来源仍能产生必然 HB”，也不主张新抽象域超越已有框架。区别在于完整交替观测与写覆盖等\emph{具体条件}，以及从这些条件推出全候选 RF、CO 和覆盖边参数界的证明。CAAT~\cite{haas2022} 将广泛公理一致性组合进 SMT 理论求解；本文的固定原始变量支持合同范围更窄，也未解决任意外部公式合取后的传播完备。

\paragraph{自动化内存模型与有界历史。}
Kater~\cite{kater2023} 通过自动机方法自动化弱内存元理论和一致性检查；KATI~\cite{kati2024} 进一步处理带关系交的元理论。本文沿用关系语言及禁止游走的思路，不将自动机、旋转或通用图商视为创新。Ax2Op~\cite{godbole2022} 在明确条件下将公理模型转为有界历史的操作模型。本文已知完整有限事件和候选后缀，处理顺序是编译布局而非未知未来程序执行；其额外义务是保留历史原始身份的条件支持。这是输入和输出合同的区别，不是对整个有界历史框架的严格分离。

\paragraph{RA 一致性与弱内存验证。}
Tunç 等~\cite{tunc2023} 研究给定 RF 后是否存在合法修改序的高效算法。未知 CO 本身并非本文首次处理的问题；本文的 RF 也未知，而且 CO 方向仍供用户假设和查询。Govind 等~\cite{govind2026} 给出单写者可解与两写者困难的边界。本文标志确为单写者，不能把这部分单独当作新可解类；但数据地址有两写者，P1--P3 提供额外结构限制，得到预编译全部支持的参数化结果，不与一般两写者困难性冲突。

基于 potentials 的验证~\cite{bargmann2024} 通过统一语义域和证明规则推理多种内存模型，关注程序断言与演绎证明。本文关注固定观测的来源身份及修改序支持，不替代这种程序级验证。两者都需要准确表达线程从通信中获得什么知识，但其证明对象和输出不同。

\paragraph{逻辑结构与知识编译。}
Conrado 和 Pavlogiannis~\cite{conrado2026} 提出弱内存验证与稳健性的 MSO 框架，以 RF 图表达并抽象修改序，适用的模型范围更广。本文保留原始 CO 查询位，重点证明\emph{选定 RF 之前}的候选并图可被可靠局部化。得到稀疏表示后，MSO／有界树宽方法同样可能适用，故不能声称超出所有结构方法。Kučera 和 Savický~\cite{kucera2021} 系统讨论知识编译到传播编码，并区分原始域一致性与含辅助变量的传播完备。本文采用前一种合同，通用支持编码不构成新的核心定理。

\begin{table}[tb]
\caption{与最近路线逐项对照。表中“本文增加”指本稿需独立完成的证明义务，不声称相应框架无法扩展。}\label{tab:related}
\small\begin{tabularx}{\textwidth}{p{2.4cm}Y Y}\toprule
路线 & 已有假设与输出 & 本文增加的具体义务\\\midrule
关系分析~\cite{haas2023} & 一般 may/must 过滤、保语义 SMT 简化 & P1--P3 蕴含全候选跨度；可接入而非框架分离\\
RF 一致性~\cite{tunc2023} & RF 已定，寻找合法修改序 & RF 仍未定，CO 身份保留，任意假设的全支持\\
单／多写者~\cite{govind2026} & 按每地址写者数划分判定复杂度 & 允许数据两写者，额外周期／密度条件控制接口\\
MSO 结构方法~\cite{conrado2026} & RF 图、逻辑表达、结构验证／稳健性 & 选取 RF 前局部化候选并图，保留 CO 位\\
知识编译~\cite{kucera2021} & 已有编译表示到传播编码 & 前端保证与身份标签使支持合同适用于本片段\\\bottomrule
\end{tabularx}
\end{table}

\section{局限与结论}\label{sec:conclusion}
本文的结构定理依赖明示而较强的程序条件。完整交替观测用于锁定来源块；阶段密度将跨度转为接口；写覆盖约束固定 CO 的生成边。删除这些条件不能沿用现有证明。网络直径和每轮负载若随实例增长，也必须计入参数。语义限定为固定事件 RA，尚不包括动态控制流、RMW、一般屏障或未知未来。将定理扩展到这些范围需要新的推理，而不是直接改换实现选项。

在这些条件内，本文证明了一个可检验的联系：来源身份未决并不妨碍获得共同因果分隔；共同因果分隔又能约束\emph{整个}候选表示，而不是只约束已选执行。身份敏感 CO 查询和完整日志修复说明了为什么公共知识与具体身份必须同时保留。原生分析对照把共同前序规则放回已有过滤框架，核心贡献因此集中在具体条件、参数界及其支持用途。

实验验证了这条联系的可执行性，也显示其实际成本边界。版本交换族能承载 32 轮完整诊断及独立选定支持核对，三线程环改变拓扑后仍满足结构解释；恒定载荷六轮与环四轮的失败则说明结构上有界不等于工程上廉价。后续研究应关注更宽松的观测条件、减少联合状态常数以及真实日志接口。当前结果提供的是一条有证明、有任务、有可复现边界的精确弱内存支持分析路线。
